\subsection{Architecture and Interface}

The simulation is architected as a custom environment compliant with the Gymnasium (formerly
OpenAI Gym) interface~\cite{towers2024gymnasium}, facilitating seamless integration with
standard DRL libraries such as Stable-Baselines3~\cite{raffin2021stable}. The environment encapsulates the complex cyber-physical interactions into a
discrete time-stepping logic.

The core logic resides in the \lstinline{step(action)} function, which executes the following
sequence at every iteration:
\begin{enumerate}
    \item \emph{Physics Update}: The UAV's position \(p_t\) is updated based on the selected
    motion primitive, and battery energy \(E_t\) is decremented according to the power
    consumption model.
    \item \emph{Channel Propagation}: The Euclidean distances to all sensors are recalculated.
    The Two-Ray Ground Reflection model is applied to determine the instantaneous RSSI for
    every ground node.
    \item \emph{Protocol Logic}: The internal states of the sensors are updated. This includes
    incrementing the Age of Information (AoI) counters, generating new data packets based on a
    Poisson process, and updating the moving average of RSSI values to simulate ADR latency.
    \item \emph{Collision Resolution}: The environment checks for packet collisions. If multiple
    sensors transmit on the same Spreading Factor (SF), the Capture Effect logic determines if
    the strongest signal (closest sensor) exceeds the interference threshold.
    \item \emph{Reward Calculation}: The multi-objective reward \(R_t\) is computed based on
    the outcome of the data collection attempts and the UAV's energy usage.
\end{enumerate}
This modular design ensures that the DRL agent is trained against a rigorous model of the
physical and protocol layers, rather than a simplified abstraction.

\subsubsection{Training Algorithm}

The complete training procedure, which integrates the physics-based updates (Two-Ray model)
and protocol constraints (ADR latency) with the DRL agent's decision process, is detailed in
\autoref{alg:dqn}.

\begin{algorithm}[h]
\caption{Protocol-Aware DQN Training for UAV Data Collection}
\label{alg:dqn}
\begin{algorithmic}[1]
    \State \textbf{Initialise:} UAV position $p_0$, Replay Buffer $\mathcal{D}$
    \State \textbf{Initialise:} Q-Network $Q(s, a \mid \theta)$, Target Network $Q'(s, a \mid \theta^{-})$
    \State \textbf{Initialise:} Hyperparameters: $\epsilon_{\max}, \epsilon_{\min}, \gamma, \alpha, B, C$
    \State \textbf{Initialise:} Sensor states (Buffers, AoI, SF), $\epsilon \leftarrow \epsilon_{\max}$
    \For{$\text{episode} \in [1, M]$}
        \State Reset environment: $s_0 \leftarrow (p_{\text{start}}, E_{\max}, \ldots)$
        \State $t \leftarrow 0$
        \While{$E > 0$ \textbf{and} Data Remaining}
            \If{$\text{random()} < \epsilon$}
                \State $a_t \leftarrow \text{RandomAction}$ \Comment{Explore}
            \Else
                \State $a_t \leftarrow \arg\max_{a}\, Q(s_t, a; \theta)$ \Comment{Exploit}
            \EndIf
            \State $s_{t+1},\, r_t,\, \text{done} \leftarrow \text{Environment.Step}(a_t)$
            \State Store $(s_t, a_t, r_t, s_{t+1})$ in $\mathcal{D}$
            \If{$|\mathcal{D}| \geq B$}
                \State Sample random minibatch $\{(s_j, a_j, r_j, s'_j)\}_{j=1}^{B}$ from $\mathcal{D}$
                \State Compute TD target: $y_j \leftarrow r_j + \gamma \max_{a'} Q'(s'_j, a' \mid \theta^{-})$
                \State Compute loss: $L(\theta) \leftarrow \frac{1}{B}\sum_{j=1}^{B}(y_j - Q(s_j, a_j \mid \theta))^2$
                \State Perform gradient descent: $\theta \leftarrow \theta - \alpha \nabla_\theta L(\theta)$
            \EndIf
            \State $s_t \leftarrow s_{t+1}$, $t \leftarrow t + 1$
        \EndWhile
        \State $\epsilon \leftarrow \max\!\left(\epsilon_{\min},\ \epsilon - \frac{\epsilon_{\max} - \epsilon_{\min}}{M}\right)$
        \If{$t \bmod C = 0$}
            \State $\theta^{-} \leftarrow \theta$
        \EndIf
    \EndFor
\end{algorithmic}
\end{algorithm}

In the consolidated format (\autoref{alg:dqn}), the environment's complexity is contained
within the \lstinline{Environment.Step($a_t$)} call. This single operation sequentially
performs the kinematic update, computes the RSSI via the Two-Ray Ground Model, executes the
ADR protocol logic, determines packet capture based on the Capture Effect, and calculates the
final multi-objective reward signal before transitioning the system to state \(s_{t+1}\).

\subsubsection{Hyperparameters}

The hyperparameters used for training the DQN agent were selected based on preliminary tuning
experiments to balance sample efficiency with stability. The specific values are detailed in
\autoref{tab:hyperparams}.

\begin{table}[h]
    \centering
    \caption{DQN Hyperparameters and Training Configuration}
    \label{tab:hyperparams}
    \begin{tabular}{ll}
        \toprule
        \textbf{Parameter}               & \textbf{Value}             \\
        \midrule
        Total Timesteps                  & \(1 \times 10^6\)          \\
        Learning Rate (\(\alpha\))       & \(1 \times 10^{-4}\)       \\
        Discount Factor (\(\gamma\))     & \(0.99\)                   \\
        Buffer Size                      & 100{,}000 transitions      \\
        Batch Size                       & 32                         \\
        Target Update Interval (\(C\))   & 1{,}000 steps              \\
        Exploration Strategy             & \(\epsilon\)-greedy        \\
        Initial \(\epsilon\)             & 1.0                        \\
        Final \(\epsilon\)               & 0.05                       \\
        Exploration Fraction             & 10\% of total steps        \\
        Network Architecture             & MLP [512, 512, 256]        \\
        Activation Function              & ReLU                       \\
        \bottomrule
    \end{tabular}
\end{table}
